% ============================================================
% BME VIK TDK abstract ("tartalmi összefoglaló") template
% Matches the official one-page sample:
% https://portal.vik.bme.hu/files/00000047.pdf
%
% Hard rules from the official sample (do not violate):
%   - max 1 page, min ~20 lines of body text
%   - NO header, footer, or page numbers
%   - figures/equations only if strictly necessary
%   - official language is Hungarian
%   - own contribution must be clearly separated from cited work
%   - body must cover: methodology, key results, conclusions
% ============================================================

\documentclass[12pt]{article}

\usepackage[utf8]{inputenc}
\usepackage[T1]{fontenc}
\usepackage[magyar]{babel}
\usepackage[a4paper, margin=2.5cm]{geometry}
\usepackage{mathptmx}      % Times-like font
\usepackage{hyperref}
\usepackage{enumitem}

\pagestyle{empty}          % <- enforces "no header/footer/page numbers"
\setlength{\parindent}{0pt}
\setlength{\parskip}{0.6em}

\begin{document}

% ---------- Title ----------
\begin{center}
    {\Large\bfseries Hiányzó absztraktok rekonstrukciója idézési gráfokból: \\
    grounding-vizsgálat gráf-kondicionált nyelvi modellekkel}
\end{center}

\vspace{0.8em}

% ---------- Authors ----------
% Format: Name, Year+Major abbreviation, email
\begin{center}
    [Vezetéknév Keresztnév] [Évfolyam] [Szak rövidítés], [email@example.com] \\
    Unyi Dániel [Évfolyam] [Szak rövidítés], [unyi.email@example.com]
\end{center}

\vspace{0.4em}

% ---------- Consultant ----------
\begin{center}
    Konzulens: Gyires-Tóth Bálint, [Tanszék neve], [email@example.com]
\end{center}

\vspace{1.2em}

% ==========================================================
% BODY — min. 20 lines, max 1 page. Structure per the official
% sample:
%   1) Context / problem — why this matters, situate vs. prior work
%   2) What is YOUR contribution specifically (own work must be
%      unambiguous — do not let it blur into the related-work summary)
%   3) Methodology used (analytic / numerical / measurement — here:
%      dataset, model, training setup)
%   4) Key results and conclusions
%   5) (optional) core references, briefly
% Fill in / replace the placeholder paragraphs below.
% ==========================================================

[TODO: 1--2 mondat a problémáról — hiányzó/rekonstruálandó
csomópont-szövegek gráf-strukturált adatokban, és miért nem oldja
ezt meg a jelenlegi GNN+LLM vagy tisztán szöveges megközelítés.]

[TODO: saját hozzájárulás egyértelmű megfogalmazása — mit csináltunk
mi, elhatárolva a máshol elért eredményektől. Legalább 1 bekezdés,
ami az olvasó számára világossá teszi, hogy ez itt a saját munka.]

[TODO: módszertan — ogbn-arxiv-TA adathalmaz, GNN-enkódolt szomszédsági
reprezentáció mint soft prefix, LoRA-finomhangolt nyelvi modell,
kiértékelés (ROUGE / BERTScore / LLM-as-judge), ablációk (szöveg-only
vs. gráf-kondicionált, shuffled-neighbour).]

[TODO: legfontosabb eredmények és következtetések — konkrét számok,
mit mutat az abláció, mit jelent ez a "figyel-e a modell a gráfra"
kérdésre nézve.]

\vspace{1em}

% ---------- References (optional in the abstract) ----------
\textbf{Irodalom:}
\begin{enumerate}[leftmargin=1.5em]
    \item [Szerző]: [Cím]. [Folyóirat/konferencia], [évszám].
    \item [Szerző]: [Cím]. [Folyóirat/konferencia], [évszám].
\end{enumerate}

\end{document}
